Desde 2025 ha emergido una corriente de trabajo académico que plantea el
uso de modelos de lenguaje grandes (\textit{Large Language Models},
\acs{llm}) para asistir o automatizar la corrección de respuestas
manuscritas. Pensieve Grader \cite{yang2025pensieve}, operado en más de
veinte instituciones, transcribe respuestas y propone calificaciones
alineadas con la rúbrica conservando siempre una etapa de revisión humana.
Otras propuestas, como VEHME \cite{nguyen2025vehme} o el sistema de la
Universidad de Liubliana \cite{pers2026grading}, aplican modelos
multimodales al ámbito de matemáticas e ingeniería. Un denominador común
recorre estos trabajos: aunque el componente de inteligencia artificial es
sustancial, todos mantienen al humano en el bucle y ninguno sostiene que
la calificación final pueda emitirse sin revisión del docente. Esta
postura coincide con la del SCDEE, que asiste al corrector pero no lo
sustituye.
